\documentclass[openany,fontset=fandol]{ctexbook}
\usepackage{amsmath,amssymb}
\usepackage[unicode,hidelinks]{hyperref}
\setlength{\emergencystretch}{2em}
\newcommand{\note}[1]{\unskip\nobreak\footnote{译注：#1}}
\begin{document}
\title{代数几何基础第五卷预备笔记\\第16节：几何奇异零点}
\author{据用户上传英译编辑稿译出}
\date{2026年11月19日}
\maketitle
\chapter*{范围与校订状态}
本批译出底本第62页定义3及注4(a)、(b)、(c)的完整段落。下一接续点为同页注4(d)：“Under the conditions of Definition 1”。注4(d)跨页，本批尚未译入。

定义3中值域的$k$据基变换语境校为$k'$，基变换对象与截面的记号统一，并以脚注保留原貌。注4(a)仅报告作者对后续判据的引用，不冒充本批已经证明；特征2下二次型退化的约定及旧批次待核项仍保留。此为独立续译，未并入第十六版主稿。
\chapter*{§16 新 EGA V（续）}
\enlargethispage{3\baselineskip}
\paragraph{定义3.}
设$X$为域$k$上的光滑预概形，$\varphi$为$\mathcal O_X$的一个截面，且$x\in V(\varphi)$。如果对于$k$的每一个扩域$k'$，以及$X$的每一个以$k'$为值域、位于$x$的点$\xi$，在$X_{k'}$中对应的点$x'$都是$\varphi_{k'}$的奇异零点（分别为超奇异零点），则称$x$为$\varphi$相对于$k$的一个\emph{几何奇异零点}（分别为\emph{几何超奇异零点}）。\note{底本在“for every extension $k'$ of $k$”之后印作“with values in $k$”，与所述基变换及注4(a)不合；此处校为“以$k'$为值域”。其含义是一个$k$态射$\xi:\operatorname{Spec}k'\to X$，像为$x$；它诱导$X_{k'}=X\times_k\operatorname{Spec}k'$中的点$x'$。并非任意扩域都存在这样的$\xi$；量词只要求检验实际存在的点。底本写$X'_k$和$\varphi'_k$，这里统一以$X_{k'}$、$\varphi_{k'}$表示基变换及拉回截面，避免撇号位置混淆。}

\paragraph{注4.}
\begin{enumerate}
\item[(a)] 由后面将建立的判据可知，在定义3中，只须用一个以某个$k'$为值域的点作检验即可。例如，可以取$k'=k(x)$或$\overline{k(x)}$，并取以这样的$k'$为值域的典范点。\note{底本第二个$k(x)$上有横线，此处明确为余域的代数闭包$\overline{k(x)}$。典范点来自$\operatorname{Spec}k(x)\to X$，或其沿$k(x)\hookrightarrow\overline{k(x)}$的复合。作者以尚待展开的判据为依据；本批未补造该判据及其证明。}
\item[(b)] 由注2可知，$x$是$\varphi$的几何非奇异零点，当且仅当$\varphi_x\ne0$，且$V(\varphi)$在$x$处光滑。\note{这里仍在定义3的域$k$上讨论，因此“光滑”指在$k$上光滑。$\varphi_x\ne0$是局部环中的芽非零，并不表示在余域中的函数值非零。$x\in V(\varphi)$时后者为零；该非零条件不能省略，否则零截面也会被误判为非奇异。}
\item[(c)] 假设预概形$X$在另一个预概形$Y$上光滑，$\varphi$为$\mathcal O_X$的一个截面，且$x\in V(\varphi)$。若$x$在纤维$X_s$上是相对于$k(s)$的奇异零点（或超奇异零点），就称$x$为相对于$Y$的奇异零点（或超奇异零点）；这里$s$是$x$在$Y$中的像。\note{纤维上所用函数是$\varphi$的限制。底本此项写“singular zero (or supersingular) relative to $Y$”和“such relative to $k(s)$”，没有在这两处重复“geometrically”；正文保留这一措辞。在紧接定义3的语境中，“相对于$k(s)$”承接域上几何版本的定义，不应悄然改成仅检验纤维局部环的未作扩域的版本。术语是否在后文统一，仍须继续对照。}
\end{enumerate}
\end{document}
